\section{Discussion}\label{sec:discussion}
\noindent\textbf{MR Coverage Complements Execution-based Coverage.}
Our results suggest that MR coverage captures a dimension of test adequacy that is not reflected by execution-based coverage alone. While statement and branch coverage indicate whether implementation code is exercised, MR coverage considers whether behavioral relations in the inferred MR space are exercised and explicitly validated. These two perspectives therefore provide complementary information for assessing existing UI component test suites.

% \vspace{1pt}
\noindent\textbf{Interpreting Behavioral Gaps.}
The distinction between \textit{Touch} and \textit{Cover} provides additional context for interpreting behavioral gaps within the inferred MR space. Low \textit{Touch} may indicate relations that are not exercised by existing tests, whereas high \textit{Touch} but low \textit{Cover} suggests that the corresponding behaviors are exercised without explicit relation-relevant validation. Such distinctions may help prioritize where additional behavioral assertions are likely to be beneficial.

% \vspace{1pt}
\noindent\textbf{Using Inferred MRs as a Behavioral Reference.}
The inferred MR space is not intended to serve as a complete behavioral specification. Instead, it provides an empirical behavioral reference for aggregate adequacy assessment. The manual validation, issue-mapping analysis, oracle-strengthening study, and MR-relevant injected-fault study together provide supporting evidence that the inferred MR space captures many practically relevant behavioral relations, while also leaving behaviors outside the current taxonomy and inference scope.